% 完整翻译 source-reference/GUIDE.tex, 源行 1--126.
\chapter*{第三卷数学指南与目录}
\addcontentsline{toc}{chapter}{第三卷数学指南与目录}
\markboth{第三卷数学指南}{第三卷数学指南}
\thispagestyle{empty}
\begin{center}
  \vspace*{14mm}
  {\Large 目录}
\end{center}
\vspace{16mm}
\begingroup
\small
\GuideChapter{\textbf{XIX. 约化群——一般性质}, M.~Demazure 著}{1}
\GuideSection{1. 代数闭域上群的回顾}{1}
\GuideSection{2. 约化群概形. 定义与初步性质}{8}
\GuideSection{3. 约化群概形的根与根系}{12}
\GuideSection{4. 根与向量群概形}{15}
\GuideSection{5. 一个有启发性的例子}{20}
\GuideSection{6. 极大环面的局部存在性. Weyl 群}{23}
\GuideSection{参考文献}{25}
\GuideChapter{\textbf{XX. 半单秩为 \(1\) 的约化群}, M.~Demazure 著}{27}
\GuideSection{1. 初等系统. 群 \(U_{\alpha}\) 与 \(U_{-\alpha}\)}{27}
\GuideSection{2. 初等系统的结构}{35}
\GuideSection{3. Weyl 群}{44}
\GuideSection{4. 同构定理}{51}
\GuideSection{5. 初等系统的例子及其应用}{53}
\GuideSection{6. 初等系统的生成元与关系}{58}
\GuideChapter{\textbf{XXI. 根数据}, M.~Demazure 著}{63}
\GuideSection{1. 一般性质}{63}
\GuideSection{2. 两个根之间的关系}{67}
\GuideSection{3. 单根与正根}{71}
\GuideSection{4. 半单秩为 \(2\) 的约化根数据}{84}
\GuideSection{5. Weyl 群: 生成元与关系}{86}
\GuideSection{6. 根数据的态射}{90}
\GuideSection{7. 结构}{101}
\GuideSection{参考文献}{108}
\GuideChapterSplit{XXII. 约化群: 分裂结构、子群、商群}{M.~Demazure 著}{109}
\GuideSection{1. 根与余根. 分裂群与根数据}{109}
\endgroup
\clearpage
\begingroup
\small
\GuideSection{2. 分裂结构的存在性. 约化群的型}{114}
\GuideSection{3. Weyl 群}{116}
\GuideSection{4. 分裂群的同态}{118}
\GuideSection{5. \textup{(R)} 型子群}{128}
\GuideSection{6. 导出群}{168}
\GuideSection{参考文献}{175}
\GuideChapter{\textbf{XXIII. 约化群: 钉扎群的唯一性}, M.~Demazure 著}{177}
\GuideSection{1. 钉扎}{177}
\GuideSection{2. 钉扎群的生成元与关系}{182}
\GuideSection{3. 半单秩为 \(2\) 的群}{190}
\GuideSection{4. 钉扎群的唯一性: 基本定理}{202}
\GuideSection{5. 基本定理的推论}{206}
\GuideSection{6. Chevalley 系统}{209}
\GuideSection{参考文献}{213}
\GuideChapter{\textbf{XXIV. 约化群的自同构}, M.~Demazure 著}{215}
\GuideSection{1. 约化群的自同构概形}{216}
\GuideSection{2. 自同构与子群}{224}
\GuideSection{3. 约化群的 Dynkin 概形. 拟分裂群}{228}
\GuideSection{4. 约化群及约化群下主丛的等平凡性}{237}
\GuideSection{5. 伴随群或单连通群的典范分解}{243}
\GuideSection{6. 约化群的 Borel 子群的自同构}{248}
\GuideSection{7. 当 \(G\) 约化时, 函子
  \(\operatorname{Hom}_{S\text{-grp.}}(G,H)\) 的可表性}{250}
\GuideSection{8. 附录: Hensel 环上光滑群的上同调. 上同调与函子 \(\prod\)}{261}
\GuideSection{参考文献}{266}
\GuideChapter{\textbf{XXV. 存在性定理}, M.~Demazure 著}{267}
\GuideSection{1. 定理的陈述}{267}
\GuideSection{2. 存在性定理: 群块的构造}{269}
\GuideSection{3. 存在性定理: 证明的结束}{274}
\GuideSection{4. 附录}{276}
\GuideSection{参考文献}{277}
\GuideChapter{\textbf{XXVI. 约化群的抛物子群}, M.~Demazure 著}{279}
\GuideSection{1. 回顾. Levi 子群}{279}
\GuideSection{2. 抛物子群的幺幂根基的结构}{285}
\GuideSection{3. 约化群的抛物子群概形}{289}
\GuideSection{4. 两个抛物子群的相对位置}{296}
\GuideSection{5. 共轭定理}{311}
\GuideSection{6. 抛物子群与分裂环面}{320}
\GuideSection{7. 相对根数据}{322}
\endgroup
\clearpage
\begingroup
\small
\GuideSection{参考文献}{330}
\endgroup
\clearpage
\thispagestyle{empty}
\null
